\chapter*{Introduction}
\label{I.0}

We present here, in revised and completed form, a photo-offset re-edition
of the second Seminar of Algebraic Geometry of the Institut des Hautes
\'Etudes Scientifiques held in 1962 (mimeographed).

The reader is referred to the Introduction to the first of these Seminars
(cited \SGA~1 below) for the aims pursued by these seminars, and their
relations with the \'El\'ements de G\'eom\'etrie Alg\'ebrique.

The text of Expos\'es~\Ref{I} to~\Ref{XI} was written as we went along,
from my oral expos\'es and manuscript notes, by a group of auditors,
comprising I\ptbl Giorgiutti, J\ptbl Giraud, Mlle~M\ptbl Jaffe (who became
Mme~M\ptbl Hakim) and A\ptbl Laudal. These notes were originally regarded
as provisional and of very limited circulation, pending their absorption
into the EGA (an absorption which has now become at the very least
problematic, as for the other parts of the SGA). As was said in the
notice to the original edition, this ``confidential'' character of the
notes was to excuse certain ``weaknesses of style'', doubtless more
apparent in the present Seminar \SGA~2 than in the others. I have tried
as far as possible to remedy this in the present re-edition, by a
relatively close revision of the initial text. In particular I have
harmonized the numbering systems of the statements used in the different
expos\'es, by introducing everywhere the same decimal system, already
used in most of the original expos\'es of the present \SGA~2, as well as
in all the other parts of the SGA. This led me in particular to revise
entirely the numbering\nde{one has preserved the original numbering as
far as possible, adding the adverb bis when ambiguous duplicates occurred
here and there.} of the statements of Expos\'es~\Ref{III} to~\Ref{VIII}
(and consequently, of the references to those expos\'es)\sfootnote{It
goes without saying that all references to \SGA~2 appearing in the parts
of the SGA published in the Series in Pure Mathematics will refer to the
present volume, and not to the original edition of \SGA~2!}. I have also
tried to extirpate from the original text the principal typing or syntax
errors (which were numerous and troublesome). Moreover, Mme~M\ptbl Hakim
kindly undertook to rewrite Expos\'e~\Ref{IV} in a less telegraphic style
than the original expos\'e. As in the other re-editions of the SGA, I have
also added a certain number of footnotes, either to give additional
references, or to indicate the state of a question for which progress has
been made since the writing of the original text. Finally, this Seminar
has been augmented by a new expos\'e, namely Expos\'e~\Ref{XIV}, written
by Mme Mich\`ele Raynaud in 1967, which takes up and completes suggestions
contained in the ``Comments on Expos\'e~\Ref{XIII}''
(\Ref{XIII}~\Ref{XIII.6}) (written in March~1963). This expos\'e takes up
the Lefschetz-type theorems from the point of view of \'etale cohomology,
using the results on \'etale cohomology set out in \SGA~4 and \SGA~5 (to
appear in this same collection Series in Pure Mathematics)\nde{in fact
these seminars are published by Springer (numbers 269, 279, 305 and~589),
but, alas, are out of print.}; it is therefore in this respect of a less
``elementary'' nature than the other expos\'es of the present volume,
which use little more than the substance of Chapters~I to~III of EGA.
Here is a sketch of the contents of the present volume. Expos\'e~\Ref{I}
contains the sorites of ``cohomology with support in~$Y$''
$H^*_Y(X, F)$, where $Y$ is a closed subset of a space~$X$, a cohomology
which can be interpreted as a cohomology of~$X$ \emph{modulo} the open
set $X-Y$, and which is the abutment of a most useful ``\emph{spectral
sequence of passage from local to global}'' \Ref{I}~\Ref{I.2.6},
involving \emph{sheaves} of cohomology ``with support in~$Y$''
$\sheaf{H}^*_Y(F)$\nde{one has retained the underlined notations for the
sheafified versions of functors, the calligraphic analogue of $\Gamma$
not being clear.}. This formalism can in many questions play a role of
``localization'' analogous to that played by the consideration of
``tubular'' neighbourhoods of~$Y$ in differential geometry.
Expos\'e~\Ref{II} studies the preceding notions in the case of
quasi-coherent sheaves on preschemes; Expos\'e~\Ref{III} gives their
relation with the classical notion of \emph{depth}
(\Ref{III}~\Ref{III.3.3}).

Expos\'es~\Ref{IV} and~\Ref{V} give notions of \emph{local duality},
which one may compare with Serre's projective duality theorem
(\Ref{XII}~\Ref{XII.1.1}); we note that these two types of duality
theorems are generalized in a substantial way in Hartshorne's seminar
(cited in a footnote at the end of \Exp~\Ref{IV})\nde{\label{noteconrad}Hartshorne's
book contains sign errors and, above all, does not really prove the
compatibility of the trace with change of base. Conrad has taken this
work up completely, proving this compatibility, which is crucial and
highly non-trivial (\textit{Grothendieck duality and base change},
Lect.\ Notes in Math., vol.~1750, Springer-Verlag, Berlin, 2000).
Alas, errors remain (\cf two preprints (Conrad~B., ``Clarifications
and corrections to ``\textit{Grothendieck duality and base change}''''
and ``An addendum to Chapter~5 of ``\textit{Grothendieck duality and
base change}'''')). For a more concrete aspect, with particular
attention to the notion of residue, see the works of Lipman, in
particular (Lipman~J., \textit{Dualizing sheaves, differentials and
residues on algebraic varieties}, Ast\'erisque, vol.~117, Soci\'et\'e
math\'ematique de France, 1984). A categorical proof of the duality
theorem, based on Brown's representability theorem, was obtained by
Neeman (Neeman~A., ``The Grothendieck duality theorem via Bousfield's
techniques and Brown representability'', \emph{J.~Amer.\ Math.\ Soc.}
\textbf{9} (1996), no.~1, p\ptbl 205--236).}

Expos\'es~\Ref{VI} and~\Ref{VII} give easy technical notions, used in
Expos\'e~\Ref{VIII} to prove the \emph{finiteness theorem}
(\Ref{VIII}~\Ref{VIII.2.3}), giving necessary and sufficient conditions,
for a coherent sheaf $F$ on a noetherian scheme~$X$, in order that the
local cohomology sheaves $\sheaf{H}^i_Y(F)$ be coherent for $i\leq n$
(or equivalently, that the sheaves $\R^if_*(F|X-Y)$ be coherent for
$i\leq n-1$, where $f\colon X-Y\to X$ is the inclusion). This theorem is
one of the central technical results of the Seminar, and we show in
Expos\'e~\Ref{IX} how a theorem of this nature can be used to establish
a ``comparison theorem'' and an ``existence theorem'' in formal
geometry, following and generalizing the use made in
(\EGA III~\oldS\S 4 and~5) of the finiteness theorem for a proper
morphism.

One applies these last results in~\Ref{X} and~\Ref{XI}, devoted
respectively to Lefschetz-type theorems for the fundamental group, and
for the Picard group. These theorems consist in comparing, under certain
conditions, the invariants ($\pi_1$ or $\Pic$) attached respectively to
a scheme~$X$ and to a subscheme~$Y$ (playing the role of a hyperplane
section), and in particular in giving conditions under which they are
isomorphic. Roughly speaking, the hypotheses made serve to pass from~$Y$
to the formal completion of~$X$ along~$Y$, and then to be able to apply
the results of~\Ref{IX} to pass from there to an open \emph{neighbourhood}
$U$ of~$Y$ in~$X$. To pass from~$U$ to~$X$, one still needs information
(of ``purity'' or ``parafactoriality'' type) for the local rings of~$X$
at the points of $Z=X-U$ (which is a finite discrete set in the cases
envisaged). This explains the interaction, in the proofs of
Expos\'es~\Ref{X}, \Ref{XI}, \Ref{XII}, between local and global
results, notably in certain inductions. The principal results obtained
in~\Ref{X} and~\Ref{XI} are the theorems \emph{of local nature}
\Ref{X}~\Ref{X.3.4} (\emph{purity theorem}) and
\Ref{XI}~\Ref{XI.3.14} (\emph{parafactoriality theorem}). One will note
that these theorems are proved by cohomological techniques, of an
essentially global nature. In~\Ref{XII}, using the preceding local
results, one obtains the global variants of these results for projective
schemes over a field, or more generally over a more or less arbitrary
base scheme; among the typical statements, we mention
\Ref{XII}~\Ref{XII.3.5} and \Ref{XII}~\Ref{XII.3.7}.

In~\Ref{XIII}, we review some of the numerous problems and conjectures
suggested by the results and methods of the Seminar. The most
interesting, perhaps, concern the cohomological and homotopical
Lefschetz-type theorems for complex analytic spaces, \cf \Ref{XIII}
pages 26 and following\nde{essentially all the conjectures stated
in~\Ref{XIII} and \Ref{XIV} are now proved; see the footnotes of those
sections for references and comments.}. In the context of the \'etale
cohomology of schemes, the corresponding conjectures are proved
in~\Ref{XIV} by a duality technique which should apply equally in the
complex analytic case (\cf comments \Ref{XIII} p\ptbl25 and
\Ref{XIV}~\Ref{XIV.6.4}). But the corresponding homotopical statements
in the case of analytic spaces (and more particularly the statements
involving the fundamental group) seem to require entirely new techniques
(\cf \Ref{XIV}~\Ref{XIV.6.4}).

I am happy to thank all those who, in various capacities, have helped
with the appearance of the present volume, including the collaborators
already cited in this Introduction. More particularly, I wish to thank
Mlle Chardon for the good grace with which she discharged the thankless
task constituted by the material preparation of the final manuscript for
photo-offset.

\bigskip\hfill Bures-sur-Yvette, April~1968
\par\smallskip \hfill A\ptbl Grothendieck.
